% Generated draft: requires mathematical and source review.
\documentclass[11pt]{article}
\usepackage[margin=25mm]{geometry}
\usepackage{fontspec}
\setmainfont{Latin Modern Roman}
\usepackage{amsmath,amssymb,mathtools,unicode-math}
\setmathfont{Latin Modern Math}
\usepackage{enumitem,needspace,xurl,hyperref}
\setlength{\parindent}{0pt}
\setlength{\parskip}{5pt}
\setlength{\emergencystretch}{3em}
\newcommand{\N}{\mathbb{N}}
\newcommand{\Z}{\mathbb{Z}}
\newcommand{\Q}{\mathbb{Q}}
\newcommand{\R}{\mathbb{R}}
\newcommand{\C}{\mathbb{C}}
\newcommand{\F}{\mathbb{F}}
\newcommand{\A}{\mathbb{A}}
\newcommand{\PP}{\mathbb{P}}
\newcommand{\E}{\mathbb{E}}
\newcommand{\eps}{\varepsilon}
\newcommand{\dd}{\,\mathrm{d}}
\begin{document}
\section*{UnsolvedMath ID: 800663}
\textbf{Source problem code:} catalogue entry
\subsection{Definitions and mathematical statement}
The supplied excerpt leaves the underlying class of spaces and the cover conventions unspecified. The following is a parameterized formulation requiring confirmation against the source. Fix a topological space \(X\), and put \(\N=\{1,2,\ldots\}\). Let \(\mathscr C(X)\) denote the admissible open covers of \(X\). For this formulation, use the nontrivial-cover convention: each member of a cover is open, its union is \(X\), and \(X\) itself is not a member. Whether admissible covers must also be countable is unresolved; the same choice of convention must be used throughout.

For \(n\in\N\), define
\[
\mathcal O_n(X)=\left\{\mathcal U\in\mathscr C(X):\ \forall F\subseteq X\text{ with }|F|\le n,\ \exists U\in\mathcal U\ (F\subseteq U)\right\}.
\]
Define the family of admissible \(\omega\)-covers by
\[
\Omega(X)=\left\{\mathcal U\in\mathscr C(X):\ \forall F\subseteq X\text{ finite},\ \exists U\in\mathcal U\ (F\subseteq U)\right\}.
\]
Under the usual groupability interpretation, \(\Omega^{gp}(X)\) consists of those \(\mathcal U\in\Omega(X)\) admitting a partition
\[
\mathcal U=\bigsqcup_{k\in\N}\mathcal B_k
\]
into nonempty finite families such that
\[
\forall F\subseteq X\text{ finite},\ \exists k_0\in\N\ \forall k\ge k_0\ \exists U\in\mathcal B_k\quad F\subseteq U.
\]
This interpretation of the superscript \(gp\), like the cover conventions, is not defined in the supplied excerpt and requires source confirmation.

For \(\mathfrak V\in\{\Omega,\Omega^{gp}\}\), the property \(\mathsf S_{\mathrm{fin}}(\Omega,\mathfrak V)\) means that for every sequence \((\mathcal U_n)_{n\in\N}\) from \(\Omega(X)\), there are finite, possibly empty, families \(\mathcal F_n\subseteq\mathcal U_n\) such that
\[
\bigcup_{n\in\N}\mathcal F_n\in\mathfrak V(X).
\]
In the game \(\mathsf G_{\mathrm{fin}}(\{\mathcal O_n\}_{n\in\N},\mathfrak V)\), inning \(n\) consists of ONE choosing \(\mathcal U_n\in\mathcal O_n(X)\) and TWO choosing a finite, possibly empty, subfamily \(\mathcal F_n\subseteq\mathcal U_n\). TWO wins exactly when \(\bigcup_{n\in\N}\mathcal F_n\in\mathfrak V(X)\); otherwise ONE wins. A winning strategy for ONE must win against every sequence of legal responses by TWO.

The question attributed to Tsaban, Problem 36 (2003), is whether, for each \(\mathfrak V\in\{\Omega,\Omega^{gp}\}\),
\[
\text{ONE has no winning strategy in }\mathsf G_{\mathrm{fin}}(\{\mathcal O_n(X)\}_{n\in\N},\mathfrak V(X))
\quad\Longleftrightarrow\quad
\mathsf S_{\mathrm{fin}}(\Omega(X),\mathfrak V(X)).
\]
The intended quantification over spaces \(X\) is missing from the supplied record. Thus this is a statement template under explicitly stated conventions, not a verified assertion that the source asks this for all topological spaces. The supplied registry labels the problem open; that label is not an independently verified assessment.
\subsection{Short English statement}
For each of the targets \(\Omega\) and \(\Omega^{gp}\), must the finite selection property for sequences of \(\omega\)-covers be equivalent to ONE having no winning strategy when ONE must play an open \(n\)-cover in inning \(n\)? The intended class of spaces and the precise cover conventions require source confirmation.
\subsection{Sources}
\begin{itemize}
\item Local UnsolvedMath registry: supplied complete record, permanent ID \(800663\), exact source problem code \texttt{catalogue entry}, title ``Selection Principles Milestone Problems --- Problem 36.'' The record attributes the problem to Tsaban (2003) and labels it open; this status is supplied metadata, not independently verified.
\item Tsaban, \emph{Selection Principles in Mathematics: A Milestone of Open Problems} (2003), Problem 36. Supplied source URL: \url{https://arxiv.org/abs/math/0312182}. The background records extraction by \texttt{arxiv\_latex} and an access date of 2026-10-16. That access date is supplied provenance; no independent source visit or literature verification is claimed here.
\end{itemize}
\end{document}
